\documentclass[main.tex]{subfiles}
\begin{document}

\qs{Question 6}{
(4 pts) Consider a Markov reward process consisting of a ring of three states \texttt{a}, \texttt{b}, and \texttt{c}, with state transitions going deterministically around the ring in that order. A reward of $+1$ is received upon leaving \texttt{a}, a reward of $+2$ on leaving \texttt{b}, and a reward of $+3$ on leaving \texttt{c}. First compute the average reward, denoted $r(\pi)$. Then use the Bellman equations and (2) (and not (1)) to compute the differential values of the three states.
}

\sol

\end{document}
